\newif\ifglobal

%%%%%%%%%%%%%%%%%%%%%%%
%%% Compile to view %%%
%%%%%%%%%%%%%%%%%%%%%%%

%%%%%%%%%%%%%%%%%%%%%%%%%%%%%%%%%%%%%%%%%%%%%%%%%%%
%%% Readme:                                     %%%
%%% https://www.overleaf.com/read/bwdjvfjksvjy/ %%%
%%%%%%%%%%%%%%%%%%%%%%%%%%%%%%%%%%%%%%%%%%%%%%%%%%%

% >>> M?: marks a module setting number or important notice

% >>> M4: Set {./relative-path} to external files for this module <<<
\def \modulefiles {./50-USBOTGHS}
% To add an external file, use:
% (Replace '---' with actual filename)
% '\input{\modulefiles/tables/---.tex}' for tables
% '\includegraphics{\modulefiles/figures/---}' for figures
% '\subfile{\modulefiles/---}' for register files

% >>> M1: This module is in global project? <<<
% 'YES' - uncomment; 'NO' - comment out
\globaltrue

\ifglobal
    \documentclass[main\_\_CN.tex]{subfiles}
   
\else
    % Font "FandolSong-Regular" used by ctex does not have GSUB table.
    % It causes warnings that do not affect compilation and can be ignored.
    \PassOptionsToPackage{quiet}{fontspec}

    \documentclass[openany, 10pt]{book}
   
    % >>> M2: In ./00-shared/preamble-trm-module, also find
    %         settings for visibility of todo notes, watermark <<< 
    \usepackage{preamble-trm-module}

    % Enable Chinese typesetting
    \usepackage{ctex}

    % >>> M3: Temporary labels in Chinese
    %         you can add your own labels there <<<
    \myexternaldocument{./00-shared/temp-labels-cn}

\fi %\ifglobal


\setcounter{secnumdepth}{4}

% Define formatting of chapters in text
\titleformat{\chapter}
  [display]
  {\LARGE\bfseries}
  {\color{AllportsColor}第\, \thechapter\, 章}
  {1em}
  {\color{AllportsColor}}
\titlespacing*{\chapter}{0pt}{0pt}{20pt}

% Define headers
\usepackage{fancyhdr}
\usepackage{lipsum}
\renewcommand{\chaptermark}[1]{\markboth{第\, \thechapter\, 章 \  \ #1}{}}
\fancypagestyle{plain}{%
  \fancyhead{}%
  %\fancyhead[R]{\Navigation}
  \fancyhead[L]{%
    \nouppercase{%
      \ %
      \normalsize\leftmark
    }%
  }%
}
\pagestyle{plain}

% >>> M5: If this module needs any updates in
% './00-shared/preamble-trm-module.sty'
% (include additional package, etc.), contact your module PM
% or create the file './00-shared/changelog.tex' make the
% necessary updates yourself and document them in this file <<<

% Make text left-aligned and allow latex to break pages naturally, comment out for full justification
\raggedright
\raggedbottom


\begin{document}


% Sets language into which pre-defined bits of text are translated
\selectlanguage{chinese}

% Do not delete this block
\ifglobal
% Add the following front matter if
% this module is outside of global project
\else
    \InputIfFileExists{readme}{}{}
    \listoftodos
    {\let\clearpage\relax \listofdonetodos}
    \tableofcontents
    %\listoftables
    %\listoffigures
\fi




%%%%%%%%%%%%%%%%%%%%%%%%%
%%% TRM Module Begins %%%
%%%%%%%%%%%%%%%%%%%%%%%%%

\hypertarget{usbotghs}{}
\chapter{USB 2.0 高速 OTG}
\label{mod:usbotghs}

\section{概述}
\chipname{} 带有一个集成了收发器的 USB 2.0 高速 OTG 外设，下文将称为 OTG\_HS。该 OTG\_HS 外设符合 USB 2.0 协议规范，同时兼容 OTG 1.3 协议和 OTG 2.0 协议。OTG\_HS 支持 USB 2.0 传输速率为 480 Mbit/s 的高速模式 (High-Speed, HS)、传输速率为 12 Mbit/s 的全速模式 (Full-Speed, FS) 和传输速率为 1.5 Mbit/s 的低速模式 (Low-Speed, LS)。 
\begin{itemize}
    \item 处于高速模式和全速模式的 OTG\_HS 可配置成 Host，也可以配置成 Device。
    \item 处于低速模式的 OTG\_HS 只可配置成 Host。
\end{itemize}

\section{术语}

本章使用到的缩写及术语如下：

\begin{longtable}[c]{ R{5cm} p{10cm} }
\textbf{主机 (Host)}     & 主机是安装了 USB 主机控制器的计算机系统，包括：主机硬件平台（CPU、总线等）和使用的操作系统。\\
\textbf{设备 (Device)}   & 设备是执行某个功能的逻辑或物理实体。具体描述的实体取决于上下文。在最低层次上，设备可能指一个单一的硬件组件，例如存储设备。在较高层次上，设备可能指执行特定功能的一组硬件组件，例如 USB 接口设备。在更高层次上，设备可能指连接到 USB 的实体，例如数据/传真调制解调器设备。非特指时，USB 设备可以是集线器或功能设备。\\
\textbf{高速 (High-speed)}       & USB 运行在 480 Mbit/s 速率下。\\
\textbf{全速 (Full-speed)}  & USB 运行在 12 Mbit/s 速率下。\\ 
\textbf{低速 (Low-speed)} & USB 运行在 1.5 Mbit/s 速率下。\\
\textbf{端点 (Endpoint)}& USB 设备的一个可以唯一寻址的部分，是在主机和设备之间通信流中信息的源或接收点。\\
\textbf{Scatter/Gather DMA 模式} & 通过 DMA 访问不连续的存储区域。\\
\textbf{Buffer DMA 模式}& 通过 DMA 访问连续存储区域。\\
\textbf{Slave 模式}& 通过 CPU 访问存储区域。\\
\textbf{RX FIFO} & 存储接收数据。 \\
\textbf{TX FIFO} & 存储发送数据。\\
\textbf{周期性 TX FIFO} & 存储待发送的同步传输和中断传输数据。\\
\textbf{非周期性 TX FIFO} & 存储待发送的控制传输和批量传输数据。\\
\end{longtable}

\section{特性}

\subsection{通用特性}
\begin{itemize}
    \item 兼容 USB 2.0 协议、OTG 1.3 协议、OTG 2.0 协议
    \item 支持高速速率、全速速率、低速速率
    \item 在全速和高速模式下既可作为主机，也可以充当设备
    \item 动态分配 FIFO (DFIFO) 大小，每个设备 EP/主机通道最大可动态分配 4 KB FIFO    
    \item 支持多种存储器访问模式
    \begin{itemize}
        \item Scatter/Gather DMA 模式
        \item Buffer DMA 模式
        \item Slave 模式
    \end{itemize}
    \item 集成 UTMI 高速收发器
\end{itemize}

\subsection{设备模式特性}
\begin{itemize}
\item 端点 0 永远存在，双向控制，由 EP0 IN 和 EP0 OUT 组成
\item 15 个附加端点 1 \verb+~+ 15，可配置为 IN 或 OUT
\item 最多 8 个 IN 端点同时工作，包括 EP0 IN
\item 所有 OUT 端点共享一个 RX FIFO
\item 每个 IN 端点都有专用的 TX FIFO
\end{itemize}

\subsection{主机模式特性}
\begin{itemize}
\item 16 个主机通道
\item 一个 RX FIFO：由所有周期事务和非周期事务共用
\item 两个 TX FIFO：
\begin{itemize}
    \item 所有非周期事务传输共用一个 TX FIFO
    \item 所有周期事务传输共用另一个 TX FIFO
\end{itemize}
\item 上述所有 FIFO 共用 4 KB RAM
\item 每个 FIFO 大小可配置，最大 4 KB
\end{itemize}

\section{功能描述}
\subsection{控制器内核与接口}

\begin{figure}[H]
\includegraphics[width=0.8\textwidth]{\modulefiles/figures/50-USBOTGHS-arch_cn.eps}
\centering
\caption {OTG\_HS 系统架构}\label{fig:usbotghs-architec}
\end{figure}

OTG\_HS 采用与 \hyperref[mod:usbotghs]{OTG\_FS} 相同的控制器内核，称为 USB 控制器内核。如图 \ref{fig:usbotghs-architec} 所示，控制器内核有以下接口：

\begin{itemize}
    \item {\bfseries CPU 接口}\\
   CPU 可以通过该接口读写控制器内核的多个寄存器和 FIFO。该接口在内部实现为 AHB 从机接口。通过该接口访问 FIFO 的方式称为 Slave 模式。
    
    \item {\bfseries DMA 接口}\\
    控制器内核的内部 DMA 可以通过该接口读写系统存储器，例如在 DMA 模式下获取和写入有效数据。该接口在内部实现为 AHB 主机接口。
    
    \item {\bfseries USB 2.0 接口 (USB 2.0 HS PHY)}\\
   控制器内核通过该接口连接 USB 2.0 UTMI 串行收发器。该串行收发器使得 OTG\_HS 可以支持高速和全速速率。
    
    \item {\bfseries 数据 FIFO 接口}\\
    控制器内核使用的多个 FIFO 实际上并不位于控制器内核内部，而是位于单端口 RAM (SPRAM) 上。FIFO 的大小可动态配置，运行时在 SPRAM 中进行分配。CPU、DMA 或控制器通过该接口读写 FIFO。
\end{itemize}

\subsection{存储器布局}

图 \ref{fig:usbotghs-reg-layout} 展示了用于配置和控制 USB 控制器内核的寄存器布局。

\begin{figure}[H]
\includegraphics[width=0.6\textwidth]{\modulefiles/figures/50-USBOTGHS-reg-layout__CN}
\centering
\caption {内核地址映射}\label{fig:usbotghs-reg-layout}
\end{figure}

\subsubsection{控制 \& 状态寄存器 (CSR)}

\begin{itemize}
    \item {\bfseries 内核全局 CSR}\\
    内核全局寄存器用于配置/控制 OTG\_HS 的通用功能，即主机和设备模式共同的功能，并表示其内部状态。通用功能包括 OTG 控制， 以及系统级中断。在主机和设备模式下，软件均可以对内核全局 CSR 进行访问。
    
    \item {\bfseries 主机模式 CSR}\\
    主机模式寄存器用于主机模式下的配置/控制/状态表示，只能在主机模式下被访问。每个通道各自有一组主机模式寄存器。
    
    \item {\bfseries 设备模式 CSR}\\
    设备模式寄存器用于设备模式下的配置/控制/状态表示，只能在设备模式下被访问。每个端点各自有一组设备模式寄存器。
    
    \item {\bfseries 电源和时钟门控寄存器}\\
    此单一寄存器用于控制模块电源和门控时钟。
\end{itemize}

\subsubsection{FIFO 访问}

OTG\_HS 利用多个 FIFO 缓冲发送或接收有效数据。FIFO 的数量和类型取决于主机或设备模式，以及所使用的通道或端点的数量，参考章节 \ref{sec:usbotghs-fifo-queue-orga}。FIFO 访问有两种方式：DMA 模式和 Slave 模式。当使用 Slave 模式时，CPU 需要通过读写 DFIFO 的推入/弹出操作 (push/pop) 区域或读写调试区域来访问这些 FIFO。FIFO 访问遵循以下规则：

\begin{itemize}
\item 对 4 KB 推入/弹出区域中任何地址的读访问将在共享 RX FIFO 中产生一个弹出 (pop) 操作。
\item 对特定 4 KB 推入/弹出区域的写访问将写入相应的端点或通道的 TX FIFO。前提是该端点是 IN 端点，或者该通道是 OUT 通道。
\begin{itemize}
\item 在设备模式下，数据写入相应的 IN 端点的专用 TX FIFO。
\item 在主机模式下，根据通道是非周期性通道还是周期性通道，数据写入非周期性 TX FIFO 或周期性 TX FIFO。
\end{itemize}
\item 访问 128 KB 读写调试区域将直接读/写，而不是进行推入/弹出操作。此种访问通常仅用于调试目的。
\end{itemize}

请注意，仅在 Slave 模式下，才需要由 CPU 直接向 FIFO 进行数据操作。在 DMA 模式下，内部 DMA 将处理 TX FIFO 和 RX FIFO 的数据推入/弹出操作。

\subsection{FIFO 和队列组织}
\label{sec:usbotghs-fifo-queue-orga}

OTG\_HS 中的 FIFO 主要用于保存有效数据，即USB 数据包的数据字段。TX FIFO 用于存储将由主机模式下的 OUT 事务或设备模式下的 IN 事务发送的有效数据。RX FIFO 用于存储主机模式下 IN 事务或设备模式下 OUT 事务的已接收的有效数据。除了存储有效数据之外，RX FIFO 还存储每个有效数据的\textbf{状态条目}。状态条目中包含关于有效数据的信息，如：通道编号、字节数、是否有效等。在 Slave 模式下，状态条目还用于指示各类通道事件。

可用于 FIFO 分配的 SPRAM 大小为 1024 × 35 位。35 位包括 32 个数据位加 3 个控制位。在主机模式下各个通道或设备模式下端点使用的多个 FIFO 被分配到 SPRAM 中，并且可以动态调整大小。

\subsubsection{主机模式 FIFO 和队列}

如图 \ref{fig:usbotghs-fifo-host-mod} 所示，主机模式使用以下 FIFO：

\begin{figure}[H]
\includegraphics[width=0.7\textwidth]{\modulefiles/figures/50-USBOTGHS-fifo-host-mod__CN}
\centering
\caption {主机模式 FIFO}\label{fig:usbotghs-fifo-host-mod}
\end{figure}

\begin{itemize}
\item {\bfseries 非周期性 TX FIFO}：存储所有通道的批量和控制类型的 OUT 事务的有效数据。
\item {\bfseries 周期性 TX FIFO}：存储所有通道的中断或同步类型的 OUT 事务的有效数据。
\item {\bfseries RX FIFO}：存储所有 IN 事务的有效数据，以及用于指示有效数据大小和事务/通道事件的状态条目。例如传输完成或通道暂停。
\end{itemize}

除 FIFO 外，主机模式还包含两个请求队列，用于存储来自多个通道的事务请求。 请求队列中的每个条目都包含 IN/OUT 通道编号以及执行事务的其他信息，例如事务类型。请求队列也用于存储其他请求类型，例如通道暂停请求。

与 FIFO 不同，请求队列的大小是固定的，且不能由软件直接访问。相反，一旦启用了通道，主机内核会自动将请求写入请求队列。请求写入队列的顺序决定了 USB 事务的处理顺序。

主机模式包含以下请求队列：

\begin{itemize}

\item {\bfseries 非周期性请求队列}：针对非周期性批量和控制事务的请求队列，最多可以存储 4 个条目。
\item {\bfseries 周期性请求队列}：针对周期性中断和同步事务的请求队列，最多可以存储 4 个条目。
\end{itemize}

调度事务时，硬件将首先执行周期性请求队列上的所有请求，再执行非周期性请求队列上的请求。

\subsubsection{设备模式 FIFO}

如图 \ref{fig:usbotghs-fifo-device-mod} 所示，设备模式使用以下 FIFO：

\begin{figure}[H]
\includegraphics[width=0.7\textwidth]{\modulefiles/figures/50-USBOTGHS-fifo-device-mod__CN.pdf}
\centering
\caption {设备模式 FIFO}
\label{fig:usbotghs-fifo-device-mod}
\end{figure}

\begin{itemize}
\item {\bfseries RX FIFO}：存储数据包内接收的有效数据和状态条目，用于指示有效数据的大小。
\item {\bfseries 专用 TX FIFO}：每个使能的 IN 端点都有一个专用的 TX FIFO，用于存储该端点的所有 IN 有效数据，而无论事务类型，包括周期性或非周期性 IN 事务。
\end{itemize}

由于有专用的 FIFO，设备模式不使用任何请求队列。IN 事务的顺序由主机确定。

\subsection{中断层次结构}

OTG\_HS 有一条中断线，可以通过中断矩阵连接到一个 CPU。可以通过置位 USB\_GLBLINTRMSK 来显示中断信号。OTG\_HS 中断是 USB\_GINTSTS\_REG 寄存器中所有位的逻辑或 (OR)，且置位 USB\_GINTMSK\_REG 寄存器中的相应位可以使能 USB\_GINTSTS\_REG 中的位。USB\_GINTSTS\_REG 包含系统级中断，还包含主机或设备模式专有的中断位以及 OTG 有关中断。 OTG\_HS 中断源的层次结构如图 \ref{fig:usbotghs-interrupt-hierar} 所示。

\begin{figure}[H]
\includegraphics[width=1\textwidth]{\modulefiles/figures/50-USBOTGHS-interr-hierar__CN.eps}
\centering
\caption {OTG\_HS 中断层次结构图}
\label{fig:usbotghs-interrupt-hierar}
\end{figure}

USB\_GINTSTS\_REG 寄存器的以下位指示较低层级的中断源：

\begin{itemize}

\item {\bfseries USB\_PRTINT} 表示主机端口有未处理的中断。USB\_HPRT\_REG 寄存器指示中断源。
\item {\bfseries USB\_HCHINT} 表示一个或多个主机通道有未处理的中断。通过读取 USB\_HAINT\_REG 寄存器可以确定哪些通道有未处理的中断，然后查询该通道的 USB\_HCINT\regindex{n}\_REG 寄存器以确定中断源。
\item {\bfseries USB\_OEPINT} 表示一个或多个 OUT 端点有未处理的中断。通过读取 USB\_DAINT\_REG 寄存器可以确定哪些 OUT 端点有未处理的中断，然后查询 OUT 端点的 USB\_DOEPINT\regindex{n}\_REG 寄存器以确定中断源。
\item {\bfseries USB\_IEPINT} 表示一个或多个 IN 端点有未处理的中断。通过读取 USB\_DAINT\_REG 寄存器可以确定哪些 IN 端点有未处理的中断，然后查询 IN 端点的 USB\_DIEPINT\regindex{n}\_REG 寄存器以确定中断源。
\item {\bfseries USB\_OTGINT} 表示 OTG 事件已触发中断。查询 USB\_GOTGINT\_REG 寄存器以确定哪些 OTG 事件触发了中断。

\end{itemize}

\subsection{DMA 模式和 Slave 模式}

USB OTG 支持 3 种存储器访问方式：Slave 模式、Buffer DMA 模式、和Scatter/Gather DMA 模式。

\subsubsection{Slave 模式}

在 Slave 模式下，所有有效数据放入 FIFO 或从 FIFO 中取出都必须通过 CPU 进行。

\begin{itemize}
\item 使用 IN 端点或 OUT 通道发送数据包时，必须将有效数据放入相应的端点或通道的 TX FIFO 中。
\item 接收到数据包时，必须先通过读取 USB\_GRXSTSP\_REG 从 RX FIFO 中取出数据包的状态条目，以确定数据包中有效数据的长度，以字节为单位。然后必须由 CPU 通过读取 RX FIFO 中的存储区域，手动从 RX FIFO 中取出相应的字节数。
\end{itemize}

\subsubsection{Buffer DMA 模式}

Buffer DMA 模式类似于 Slave 模式，不同之处在于该模式为利用内部 DMA 将有效数据放入 FIFO 或从 FIFO 中取出。

\begin{itemize}
\item 使用 IN 端点或 OUT 通道传输数据包时，应将有效数据的存储地址写入 USB\_HCDMA\regindex{n}\_REG（主机模式）或 USB\_DIEPDMA\regindex{n}\_REG（设备模式）。启用端点或通道后，内部 DMA 会将有效数据从存储器中推送到通道或端点的 TX FIFO 中。
\item 使用 OUT 端点或 IN 通道接收数据包时，应将存储器中空 buffer 的地址写入 USB\_HCDMA\regindex{n}\_REG（主机模式）或 USB\_DOEPDMA\regindex{n}\_REG（设备模式）。启用端点或通道后，内部 DMA 将把有效数据从 RX FIFO 弹出到相应 buffer 中。
\end{itemize}

\subsubsection{Scatter/Gather DMA 模式}

\begin{figure}[H]
\includegraphics[width=0.7\textwidth]{\modulefiles/figures/50-USBOTGHS-dma-discriptor-list__EN}
\centering
\caption {Scatter/Gather DMA 链表结构}
\label{fig:usbotghs-dma-descriptor-list}
\end{figure}

在 Scatter/Gather DMA 模式下，包含有效数据的接收 可能分散在整个存储器各处。每个端点或通道都有一个连续的 DMA 描述符列表。每个描述符包含一个指向有效数据或接收 buffer 的 32 位指针和一个 32 位的 buffer 描述符 (BufferStatus Quadlet)。有效数据和接收 buffer 可以对应单个事务，即 < 1 MPS 字节，MPS：最大数据包大小，或整个传输，即 > 1 MPS 字节。该列表的实现为环形 buffer ，意味着 DMA 在遇到列表中的最后一个条目时将返回至第一个条目。

\begin{itemize}
\item 当使用 IN 端点或 OUT 通道发送传输/事务时，DMA 将从多个 buffer 收集有效数据并将其放入 TX FIFO。
\item 当使用 OUT 端点或 IN 通道接收传输/事务时，DMA 将取出来自 RX FIFO 的接收有效数据，并将它们分别存放到 DMA 列表条目所指向的多个 buffer 中。
\end{itemize}

\subsection{事务和传输级操作}

在主机或设备模式下，通信可以在事务级或传输级进行。

\subsubsection{DMA 模式下的事务和传输级操作}

在 DMA 模式下的传输级操作，只有当一个传输通道被终止时，中断才会发生。以下三种情况下发生传输通道终止：传输通道中所有要求传送的数据全部成功传送、接收到 STALL 或出现连续事物级错误，例如，3 个连续的事务级错误。在 DMA 设备模式下操作时，所有错误均由控制器内核进行处理。

在 DMA 模式下的事务级操作，传输大小是一个数据包的大小，即最大数据包大小或短数据包大小。

\subsubsection{Slave 模式下的事务和传输级操作}

在 Slave 模式下的事务级操作，一次只能处理一个事务。每组有效数据应对应一个数据包，软件必须根据 USB 接收到的握手应答，例如 ACK 或 NAK 来确定是否需要重启事务。

下表描述了 Slave 模式下进行 IN 和 OUT 事务级操作的方法。

\input{\modulefiles/tables/50-USBOTGHS-in-and-out-transaction-device-mod__CN}

在 Slave 模式下进行传输级操作时，可以在队列中一次性排入一个或多个事务级操作，达到类似于 DMA 模式下的传输级操作的效果。在同一次触发的传输中，多个事务的数据包均可以从 FIFO 中读写，这样就无需以数据包为单位触发中断。

Slave 模式下进行传输级操作的方法类似于事务级操作，不同之处在于，需要配置 USB\_HCTSIZ\regindex{n}\_REG 或 USB\_DOEPTSIZ\regindex{n}\_REG 寄存器以指定整次传输的大小和数据包个数。使能通道或端点后，应分别向 TX FIFO 或 RX FIFO 写入或读取对应于多个数据包的有效数据，假设有足够的空间或足够的数据量。

\subsection{OTG}

OTG\_HS 在内部将 ID/VBUS 悬空，因此不支持 SRP/HNP 协议，但依旧支持 DRD 双角色设备属性，即既可以作为主机，也可以充当设备。软件可以设置 USB\_GUSBCFG 中的 ForceDevMode 或者 ForceHstMode 强制设置控制器工作在主机模式下或是设备模式下而无视 ID 管脚的值。

\section{寄存器}

USB OTG 寄存器列表及详细描述受 IP 供应商保密协议 (NDA) 的约束。若需获取某一寄存器的支持信息，请通过 \href{https://www.espressif.com/zh-hans/contact-us/technical-inquiries}{技术支持} 联系乐鑫技术支持团队。

\end{document}
